\subsection{A coherent description from simple calculations}

A complete analysis should connect domain, limits, derivatives, and conclusions rather than present isolated computations.

\begin{examplebox}[An elementary polynomial analysis]
For
\[
f(x)=x^3-3x,
\]
the domain is $\mathbb R$ and the zeros are $0,\pm\sqrt3$. The derivative
\[
f'(x)=3(x-1)(x+1)
\]
shows increase on $(-\infty,-1)$ and $(1,\infty)$ and decrease on $(-1,1)$. Hence there is a local maximum at $(-1,2)$ and a local minimum at $(1,-2)$. Since
\[
f''(x)=6x,
\]
the graph changes convexity at $(0,0)$. The leading term $x^3$ determines the opposite behaviors at the two ends.
\end{examplebox}

\begin{examplebox}[A simple rational function]
For
\[
g(x)=\frac{x+1}{x-1},
\]
the domain excludes $x=1$. The limits identify the vertical asymptote $x=1$ and horizontal asymptote $y=1$. Since
\[
g'(x)=\frac{-2}{(x-1)^2}<0,
\]
the function decreases on each interval of its domain.
\end{examplebox}

\begin{summarybox}[Function analysis]
Determine the domain, useful zeros and signs, limits at boundaries and infinity, asymptotes, first-derivative sign, extrema, second-derivative sign, and then write one coherent description of the graph or model.
\end{summarybox}

The examples remain algebraically elementary so that assessment can test whether the student understands the conclusions carried by derivatives and limits.
